%==========================================================
\section{Recolha de dados}\label{sec:recolha}

\subsection{Sessões realizadas}

Foram realizadas duas sessões de investigação contextual. A Tabela~\ref{tab:sessoes} resume cada uma. As notas de campo completas estão no Apêndice~\ref{anx:notas}.

\begin{table}[!htbp]
\centering\small
\begin{tabularx}{\textwidth}{@{}p{3.1cm}LL@{}}
\toprule
 & \textbf{A0: agente de viagens} & \textbf{U0: jovem viajante} \\
\midrule
Papel & Agente de balcão, 12 anos na mesma agência. & Organizador(a) do grupo, com traços de caçador(a) de promoções. \\
Perfil & -- & 20 anos, estudante do 2.º ano, sem aeroporto na cidade; mesada e explicações; cartão de débito, sem crédito. \\
Tema & Preocupações com a IA e impacto no negócio. & Viagem de 4 noites com 2 amigos, limite de 300\,€ por pessoa. \\
Tipo de tarefa & Normal: manhã de trabalho observada ao balcão, com relatos de casos concretos. & Longa e intermitente: \textit{work walkthrough} da viagem já planeada, com U0 a mostrar e refazer os passos. \\
Local e duração & Na agência, 11\,h15--12\,h (\(\approx\)10\,min de conversa, \(\approx\)35\,min de observação). & No quarto de U0, com portátil e telemóvel; \(\approx\)1\,h45. \\
Registo & Transcrição do conteúdo necessário, com consentimento; sem fotos. & Transcrição do conteúdo necessário, com consentimento; sem capturas de ecrã. \\
Registos produzidos & AE01--AE08 (7 observados, 1 relato) & UE01--UE19 (8 observados, 11 relatos) \\
Notas de interpretação & A0-01--A0-15 & U0-01--U0-21 \\
\bottomrule
\end{tabularx}
\caption{Resumo das sessões de investigação contextual.}
\label{tab:sessoes}
\end{table}

\subsection{Como os dados foram registados}

Cada sessão deu origem a um documento com a mesma estrutura: perfil e situação da entrevista, conversa convencional, registos da entrevista contextual por ordem cronológica, fecho validado com o participante e notas da sessão de interpretação. A Tabela~\ref{tab:siglas} resume as siglas usadas.

\begin{table}[!htbp]
\centering\small
\begin{tabularx}{\textwidth}{@{}p{3cm}L@{}}
\toprule
\textbf{Sigla} & \textbf{Significado} \\
\midrule
AE01, UE01\ldots & Registo da entrevista ao agente (AE) ou ao viajante (UE) \\
A0-01, U0-01\ldots & Nota da sessão de interpretação do caso A0 ou U0 \\
O & Observado: o entrevistador viu, ou o participante mostrou ou refez o passo durante a entrevista \\
R & Relato de um caso concreto, sem demonstração \\
F/OP & Fala geral ou opinião sem caso concreto; contextualiza mas não entra nos modelos \\
OBS, INS, CUL, Q, BD & Tipos de nota: observação, \textit{insight}, influência cultural, pergunta por esclarecer, \textit{breakdown} \\
\bolt{} & \textit{Breakdown}: falha ou atrito no trabalho \\
\bottomrule
\end{tabularx}
\caption{Siglas usadas nos registos e nos modelos.}
\label{tab:siglas}
\end{table}

\paragraph{Regras de registo.} Só entra nos modelos o que foi observado ou contado como caso concreto. Frases gerais («normalmente comparo uns três sites») ficam como contexto. Os registos não têm nomes de pessoas, instituições ou cidade de origem; os amigos de U0 aparecem como «Amigo~A» e «Amigo~B». A pedido da agência, alguns pormenores sobre o software que A0 usa foram assumidos pela equipa, para não divulgar o seu funcionamento.

\subsection{Sessões de interpretação}

Cada entrevista foi interpretada pela equipa numa sessão própria, com os papéis do Capítulo~7 do livro: o entrevistador conta a sessão pela ordem em que aconteceu e desenha o modelo físico; um modelador desenha o fluxo e o cultural; outro desenha as sequências e os artefactos. Os papéis rodaram entre as duas sessões (Apêndice~\ref{anx:organizacao}), para que ninguém modelasse sozinho a sua própria entrevista.

\subsection{Principais observações}

\begin{itemize}
  \item \textbf{O custo real está espalhado e só aparece no fim.} U0 copia preços à mão de vários sites para uma folha de cálculo e só no pagamento descobre a bagagem de porão (+34\,€) \reg{UE06, UE07}. O voo mais barato exigia chegar ao aeroporto de madrugada, sem transporte \reg{UE04}.
  \item \textbf{A decisão em grupo é lenta e o preço muda entretanto.} Dois dias sem resposta e o voo subiu 19\,€ \reg{UE08}. O grupo não abre a folha partilhada e pede «print» \reg{UE17}.
  \item \textbf{O organizador acaba por ser o banco do grupo.} U0 paga o voo dos três e cobra depois; um amigo só paga no fim do mês \reg{UE15, UE16}.
  \item \textbf{A IA é usada, mas não é de confiança.} U0 pede um roteiro ao ChatGPT e verifica-o no mapa \reg{UE10}. A0 teme uma IA que «parece certa e não é» e um supervisor que acha que ela o substitui (fecho de A0).
  \item \textbf{O agente acrescenta valor onde o planeamento do cliente falha.} A0 pergunta «o que não abdicam» e mostra que o roteiro do casal é inviável \reg{AE03, AE04}, mas o cliente sai para «comparar online» \reg{AE06} e recusa a taxa de serviço porque «o site não cobra» \reg{AE07}.
  \item \textbf{O agente perde tempo a copiar dados entre sistemas} \reg{AE05} e guarda num caderno pessoal conhecimento que não está em lado nenhum \reg{AE08}.
\end{itemize}
